\section{Introduction}
Procedural memory lets an agent reuse a successful method without deriving it again. Yet success is conditional. A record may have been complete before reaching a parser; a cache may have belonged to one request; a diagnostic object may have contained only public fields. If the target changes these assumptions, a memory can accurately describe the source task while providing an incomplete guide to the new one.

Consider the status renderer in family F08. In the source context, a status value is a decimal integer, so interpolating it into HTML cannot introduce markup. The target task changes that value to caller controlled text. A small change in the input contract, with a different security consequence. Direct reuse still renders a status message, but fails the focal injection witness; a correct adaptation escapes the text before interpolation. The memory remains accurate about its source, while its applicability changes at the target boundary.

\begin{figure}[t]
\centering\includegraphics[width=\linewidth]{figures/design_overview.pdf}
\caption{Task construction and memory intervention, illustrated by F08. A procedure is valid under the source assumption. All four conditions use the same shifted target; only the supplied memory differs. Functionality and the family specific focal witness are observed separately.}
\label{fig:design}
\end{figure}

We ask: \emph{when a coding agent receives a source procedure in a target context where its assumptions no longer hold, how does that memory change functional completion and focal security failures?} We also test a generic applicability reminder. Four memory conditions share the same target task and changed assumption (Figure~\ref{fig:design}). The original hypothesis predicted that correct memory would increase functional completions that violate the focal property. That prediction is not established across the six configurations. The implementation evidence explains why the aggregate endpoint needs care: an agent can adapt the procedure, omit the required check, or avoid the violation by failing to implement the feature.

The study contributes (1) a controlled benchmark of thirteen security relevant applicability shifts; (2) a six configuration experiment that separates functionality from focal witness failure; and (3) implementation and visible trace evidence that distinguishes unchecked reuse, target appropriate protection, and failure to complete. The shifts span caller controlled paths, deferred resource lifetimes, cross user caches, origin changes, repeated control fields, concurrency invariants, and approximate authorization evidence. The comparisons preserve this task diversity while holding each target fixed across memory conditions.

All targets contain a context shift, and memories are supplied directly. The intervention therefore estimates the effect of supplied memory in those targets; it does not identify the effect of context shift or evaluate a retrieval system. Comparisons are made within each model and agent configuration.
